% 2011 年天津卷文科第 10 题：三视图（单位：m）
% 来源 PDF：普通高考/2011/2011天津文.pdf，第 1 页（印面 707）第 2 栏。
% 重绘依据：正视图总宽 2、总高 3，上块宽 1 高 2，下块宽 2 高 1；
%   侧视图宽 1、高 3，在距底边 1 处有一条水平分界线；俯视图宽 2、深 1，在 x=1 处有分界线。
\begin{tikzpicture}[scale=0.78, semithick, every node/.style={font=\small}]
  % ---------- 正视图 ----------
  \begin{scope}[shift={(0,0)}]
    \draw (0,0) rectangle (2,1);          % 下块 2×1
    \draw (0,1) rectangle (1,3);          % 上块 1×2
    \draw[<->] (-0.35,0) -- (-0.35,3) node[midway,left] {$2$};
    \draw[<->] (0,-0.35) -- (2,-0.35) node[midway,below] {$2$};
    \draw[<->] (1.35,0) -- (1.35,1) node[midway,right] {$1$};
    \node at (1,-1.0) {正视图};
  \end{scope}
  % ---------- 侧视图 ----------
  \begin{scope}[shift={(4.4,0)}]
    \draw (0,0) rectangle (1,3);
    \draw (0,1) -- (1,1);
    \draw[<->] (-0.35,0) -- (-0.35,3) node[midway,left] {$2$};
    \draw[<->] (1.3,0) -- (1.3,1) node[midway,right] {$1$};
    \draw[<->] (0,-0.35) -- (1,-0.35) node[midway,below] {$1$};
    \node at (0.5,-1.0) {侧视图};
  \end{scope}
  % ---------- 俯视图 ----------
  \begin{scope}[shift={(0,-4.6)}]
    \draw (0,0) rectangle (2,1);
    \draw (1,0) -- (1,1);
    \draw[<->] (0,-0.35) -- (2,-0.35) node[midway,below] {$2$};
    \draw[<->] (2.3,0) -- (2.3,1) node[midway,right] {$1$};
    \draw[<->] (0,1.35) -- (1,1.35) node[midway,above] {$1$};
    \node at (1,-1.0) {俯视图};
  \end{scope}
\end{tikzpicture}
